% =============================================================================
%  SECCIONES REDACTADAS — PARTE 1
%  Capítulo 1 (Introducción) y Capítulo 3 (Estado del arte)
%
%  Cada bloque indica qué sección sustituye en main.tex
% =============================================================================


% ╔═══════════════════════════════════════════════════════════════════════════╗
% ║  SUSTITUYE LA SECCIÓN 1.1                                                 ║
% ╚═══════════════════════════════════════════════════════════════════════════╝

\section{Contexto y motivación}

El sistema eléctrico español ha atravesado en los últimos años un periodo de
transformación acelerada. La progresiva integración de generación renovable, la
electrificación del transporte y la volatilidad de los mercados energéticos han
situado la gestión de la demanda en el centro del debate público y técnico.

La predicción de la demanda eléctrica constituye una función crítica para el
operador del sistema. La red carece de capacidad de almacenamiento significativa,
por lo que generación y consumo deben equilibrarse de forma continua: una
previsión deficiente obliga a recurrir a mecanismos de ajuste más costosos y
menos eficientes desde el punto de vista ambiental. Este problema ha sido
ampliamente estudiado en la literatura, con un cuerpo de trabajo consolidado en
torno a la predicción de carga a corto plazo \parencite{hong2016probabilistic}.

Existe, sin embargo, una segunda dimensión del problema que ha recibido mucha
menos atención: la del **consumidor doméstico**. En España, la tarifa regulada
—el Precio Voluntario para el Pequeño Consumidor, PVPC— establece un precio que
varía cada hora del día en función del resultado del mercado mayorista. El
consumidor acogido a esta modalidad puede, en principio, reducir su factura
desplazando el consumo hacia las horas más económicas.

En la práctica, aprovechar esa posibilidad exige una información que no está al
alcance del usuario medio. Los precios del día siguiente se publican a las 20:15
horas, una vez cerrada la subasta diaria, y su interpretación requiere una
familiaridad con el funcionamiento del mercado eléctrico que difícilmente cabe
presuponer. El resultado es una asimetría notable: la información existe, es
pública y es gratuita, pero permanece inaccesible para quien más podría
beneficiarse de ella.

Este trabajo parte de esa observación. Su motivación no es únicamente construir
un modelo predictivo preciso —objetivo ampliamente cubierto por la literatura—
sino **cerrar la distancia entre el dato abierto y la decisión del ciudadano**,
mediante un sistema que traduzca la previsión técnica en una recomendación
comprensible y accionable.


% ╔═══════════════════════════════════════════════════════════════════════════╗
% ║  SUSTITUYE LA SECCIÓN 1.2                                                 ║
% ╚═══════════════════════════════════════════════════════════════════════════╝

\section{Justificación del trabajo}

La justificación del presente trabajo se articula en torno a tres argumentos.

\subsection*{Disponibilidad de datos abiertos infrautilizados}

Red Eléctrica de España publica la demanda peninsular con resolución de diez
minutos y sin restricciones de acceso, y la Agencia Estatal de Meteorología
ofrece tanto series climatológicas históricas como previsiones a cuarenta y ocho
horas. Ambas fuentes son gratuitas y de acceso programático. Pese a ello, su
explotación conjunta con fines de asesoramiento al consumidor final resulta
escasa.

\subsection*{Distancia entre la literatura académica y la herramienta utilizable}

La investigación sobre predicción de demanda eléctrica es abundante, pero sus
resultados rara vez trascienden el entorno experimental. Un modelo evaluado en un
cuaderno de análisis y un sistema desplegado al que un usuario puede acceder son
artefactos de naturaleza distinta: el segundo exige resolver la obtención de
datos en tiempo real, el tratamiento de fallos de las fuentes externas, la
coherencia entre el código de entrenamiento y el de producción, y la traducción
del resultado numérico a lenguaje natural.

\subsection*{Aportación diferencial: un sistema completo}

El trabajo no se limita a comparar algoritmos. Abarca el ciclo completo —ingesta,
tratamiento, modelado, evaluación y despliegue— y produce como resultado una
aplicación accesible públicamente. Esta orientación condiciona decisiones
metodológicas que un estudio puramente comparativo podría obviar: la
disponibilidad efectiva de cada variable en el instante de predicción, el coste
computacional del reentrenamiento o el comportamiento del sistema ante la caída
de una fuente externa.


% ╔═══════════════════════════════════════════════════════════════════════════╗
% ║  SUSTITUYE LA SECCIÓN 1.3                                                 ║
% ╚═══════════════════════════════════════════════════════════════════════════╝

\section{Estructura del documento}

El documento se organiza en seis capítulos.

El \textbf{capítulo 2} formula el objetivo general y los siete objetivos
específicos, y delimita el alcance del trabajo.

El \textbf{capítulo 3} revisa el estado del arte: el panorama general de la
predicción de demanda eléctrica, los modelos clásicos de series temporales, las
técnicas de aprendizaje automático aplicadas a este dominio y el funcionamiento
del mercado eléctrico español.

El \textbf{capítulo 4} expone la metodología. Describe la arquitectura del
sistema, las fuentes de datos, el tratamiento aplicado —que constituye una de las
aportaciones centrales—, la ingeniería de características derivada del análisis
exploratorio, la estrategia de validación y los modelos evaluados.

El \textbf{capítulo 5} presenta y discute los resultados: la caracterización de
la serie, la comparativa de rendimiento entre las diez formulaciones evaluadas,
el análisis de importancia de variables, el estudio de errores y la aplicación
desarrollada.

El \textbf{capítulo 6} revisa el grado de cumplimiento de los objetivos, sintetiza
las aportaciones, reconoce las limitaciones y propone líneas de continuación.

Se incluyen finalmente las referencias bibliográficas y tres anexos con el
repositorio de código, la aplicación desplegada y el cronograma del proyecto.


% ╔═══════════════════════════════════════════════════════════════════════════╗
% ║  SUSTITUYE LA SECCIÓN 3.1                                                 ║
% ╚═══════════════════════════════════════════════════════════════════════════╝

\section{Predicción de demanda energética: panorama general}

La predicción de la demanda eléctrica constituye una de las aplicaciones más
antiguas y mejor documentadas del análisis de series temporales. Su relevancia
deriva de una restricción física del sistema: la electricidad no puede almacenarse
en cantidades significativas a escala de red, de modo que generación y consumo
han de equilibrarse de forma instantánea y permanente.

\subsection*{Clasificación por horizonte temporal}

La literatura distingue habitualmente cuatro horizontes, cada uno asociado a una
finalidad operativa y a metodologías distintas
\parencite{hong2016probabilistic}:

\begin{itemize}
  \item \textbf{Muy corto plazo} (de minutos a una hora). Orientado al control en
        tiempo real y a la regulación de frecuencia. Predomina la inercia de la
        serie: los métodos autorregresivos simples resultan difíciles de superar.

  \item \textbf{Corto plazo} (de una hora a una semana). Es el horizonte de la
        programación diaria y de la casación del mercado, y el que concentra el
        mayor volumen de investigación. \textbf{El presente trabajo se sitúa en
        este intervalo}, con un horizonte de veinticuatro a cuarenta y ocho horas.

  \item \textbf{Medio plazo} (de una semana a un año). Empleado en la
        planificación del mantenimiento y en la contratación de combustible.

  \item \textbf{Largo plazo} (más de un año). Orientado a la planificación de
        infraestructuras. Depende de variables macroeconómicas y demográficas más
        que de la propia serie histórica.
\end{itemize}

\subsection*{Factores determinantes}

Existe consenso en la literatura sobre los tres grupos de factores que gobiernan
la demanda a corto plazo:

\begin{enumerate}
  \item \textbf{Calendario.} Hora del día, día de la semana y festividades. Su
        efecto es sistemático y predecible, y constituye la fuente principal de
        estructura de la serie.

  \item \textbf{Meteorología.} Fundamentalmente la temperatura, cuya relación con
        la demanda no es lineal sino aproximadamente cuadrática, por la
        superposición de las demandas de calefacción y refrigeración.

  \item \textbf{Actividad económica.} Relevante en horizontes medios y largos,
        con efecto reducido a escala de días.
\end{enumerate}

\subsection*{Evolución metodológica}

El campo ha atravesado tres etapas sucesivas. Hasta los años noventa
predominaron los modelos estadísticos de la familia Box-Jenkins
\parencite{boxjenkins} y la regresión con variables meteorológicas. La década
siguiente vio la incorporación de las redes neuronales, con resultados desiguales.
Desde mediados de la década de 2010 se han impuesto los métodos de conjunto
basados en árboles de decisión.

Las competiciones M han proporcionado la evidencia comparativa más sólida al
respecto. La cuarta edición \parencite{makridakis2020m4} estableció que las
combinaciones de métodos superan sistemáticamente a los modelos individuales, y
que el aprendizaje transversal entre series mejora los resultados frente al ajuste
individualizado. La quinta \parencite{makridakis2022m5} arrojó un hallazgo aún
más rotundo: los métodos basados en árboles con potenciación del gradiente
dominaron las primeras posiciones, por delante de las arquitecturas de aprendizaje
profundo.

Esta evidencia justifica la progresión metodológica adoptada en el presente
trabajo, que parte de modelos estadísticos clásicos y escala hacia técnicas de
potenciación del gradiente, en lugar de recurrir directamente a formulaciones de
mayor complejidad.


% ╔═══════════════════════════════════════════════════════════════════════════╗
% ║  SUSTITUYE LA SECCIÓN 3.2 COMPLETA (3.2, 3.2.1, 3.2.2, 3.2.3)             ║
% ╚═══════════════════════════════════════════════════════════════════════════╝

\section{Modelos clásicos de series temporales}

Los modelos estadísticos de series temporales constituyen la referencia frente a
la que debe medirse cualquier propuesta alternativa. Su valor en este trabajo es
doble: proporcionan una línea base rigurosa y aportan herramientas de diagnóstico
—autocorrelación, descomposición estacional— que orientan la construcción de
variables en los modelos posteriores.

\subsection{Modelos ARIMA y SARIMA}

La formulación ARIMA$(p,d,q)$ \parencite{boxjenkins} modela una serie mediante
tres componentes. El término \textbf{autorregresivo} de orden $p$ expresa el valor
actual como combinación lineal de $p$ valores anteriores; la \textbf{integración}
de orden $d$ aplica diferenciaciones sucesivas hasta alcanzar la estacionariedad;
y el término de \textbf{media móvil} de orden $q$ incorpora los $q$ errores de
predicción precedentes.

Formalmente, empleando el operador de retardo $B$:

\begin{equation}
\phi_p(B)\,(1-B)^d\, y_t = \theta_q(B)\,\varepsilon_t
\end{equation}

donde $\phi_p$ y $\theta_q$ son los polinomios autorregresivo y de media móvil, y
$\varepsilon_t$ un proceso de ruido blanco.

La extensión estacional SARIMA$(p,d,q)(P,D,Q)_s$ añade componentes análogos sobre
el periodo estacional $s$:

\begin{equation}
\Phi_P(B^s)\,\phi_p(B)\,(1-B^s)^D (1-B)^d\, y_t
  = \Theta_Q(B^s)\,\theta_q(B)\,\varepsilon_t
\end{equation}

Esta formulación presenta una limitación determinante para el problema que nos
ocupa: **admite un único periodo estacional**. Una serie horaria de demanda
eléctrica exhibe al menos dos —el ciclo diario de veinticuatro horas y el semanal
de ciento sesenta y ocho—, y el coste computacional del ajuste crece con el valor
de $s$, lo que hace inviable el tratamiento directo de la estacionalidad semanal.

La solución habitual consiste en representar la estacionalidad mediante términos
de Fourier introducidos como regresores externos, reservando la parte ARIMA para
la autocorrelación de corto plazo \parencite{hyndman2021}. Este enfoque admite
varios periodos simultáneos y reduce drásticamente el coste de estimación.

\subsection{Prophet}

Prophet \parencite{taylor2018prophet} adopta una aproximación distinta:
descompone la serie de forma aditiva en tres componentes interpretables,

\begin{equation}
y(t) = g(t) + s(t) + h(t) + \varepsilon_t
\end{equation}

donde $g(t)$ representa la tendencia —lineal por tramos con puntos de cambio
detectados automáticamente—, $s(t)$ la estacionalidad modelada mediante series de
Fourier, $h(t)$ el efecto de las festividades y $\varepsilon_t$ el término de
error.

Sus ventajas para este dominio son la admisión nativa de estacionalidades
múltiples y la incorporación explícita de un calendario de festivos, con ventanas
configurables para recoger el efecto de arrastre sobre vísperas y días
posteriores. Como contrapartida, al no modelar la autocorrelación de los residuos,
tiende a desaprovechar la fuerte inercia de corto plazo característica de las
series de demanda.

\subsection{Análisis espectral y series de Fourier}

El análisis espectral permite identificar los periodos dominantes de una serie de
manera objetiva, sin presuponerlos a partir del conocimiento del dominio. La
transformada discreta de Fourier descompone la señal en componentes sinusoidales,

\begin{equation}
X_k = \sum_{n=0}^{N-1} x_n \, e^{-i 2\pi k n / N}
\end{equation}

y el periodograma $|X_k|^2$ revela la potencia asociada a cada frecuencia. Los
máximos locales identifican los ciclos presentes en los datos.

Su utilidad en este trabajo es doble. Por una parte, fundamenta empíricamente la
selección de los desfases temporales empleados como variables predictoras. Por
otra, permite representar la estacionalidad mediante pares seno-coseno en los
modelos que no la admiten de forma nativa.


% ╔═══════════════════════════════════════════════════════════════════════════╗
% ║  SUSTITUYE LA SECCIÓN 3.3 COMPLETA (3.3, 3.3.1, 3.3.2)                    ║
% ╚═══════════════════════════════════════════════════════════════════════════╝

\section{Aprendizaje automático aplicado a series temporales}

La aplicación de algoritmos de aprendizaje automático a series temporales exige
una reformulación previa del problema. Mientras los modelos estadísticos operan
directamente sobre la secuencia, los algoritmos de regresión requieren una matriz
de observaciones independientes. La transformación consiste en construir, para
cada instante, un vector de características que recoja la información disponible
hasta ese momento.

Esta reformulación introduce dos ventajas y un riesgo. Permite incorporar
variables exógenas sin restricciones y capturar relaciones no lineales e
interacciones entre predictores. Pero expone el procedimiento a la
\textbf{fuga de información}: si alguna variable incorpora datos posteriores al
instante de predicción, las métricas obtenidas serán excelentes e irreproducibles
en explotación. La sección \ref{sec:fuga} aborda esta cuestión de forma específica.

\subsection{Modelos basados en árboles}

Los métodos de conjunto basados en árboles de decisión dominan actualmente las
aplicaciones sobre datos tabulares, y las competiciones M han confirmado su
superioridad también en predicción de series temporales
\parencite{makridakis2022m5}.

\textbf{Random Forest} \parencite{breiman2001rf} construye múltiples árboles sobre
remuestras del conjunto de entrenamiento, aleatorizando además las variables
consideradas en cada división, y promedia sus predicciones. La reducción de
varianza resultante lo hace robusto frente al sobreajuste, a costa de un sesgo
algo mayor.

La \textbf{potenciación del gradiente} adopta una estrategia secuencial: cada
árbol se ajusta sobre los residuos del conjunto acumulado hasta el momento. Las
dos implementaciones de referencia son \textbf{XGBoost}
\parencite{chen2016xgboost}, que incorpora regularización explícita sobre la
complejidad del árbol y tratamiento nativo de valores ausentes, y
\textbf{LightGBM} \parencite{ke2017lightgbm}, que introduce el crecimiento por
hojas y el agrupamiento de características mutuamente excluyentes, con una
reducción sustancial del tiempo de entrenamiento.

La principal limitación de estos métodos en el contexto de series temporales es su
incapacidad para extrapolar fuera del rango observado, consecuencia directa de que
sus predicciones son promedios de valores del conjunto de entrenamiento. Esta
restricción resulta poco problemática en demanda eléctrica, cuya magnitud se
mantiene dentro de márgenes estables.

\subsection{Ingeniería de características temporales}

La calidad de las variables construidas determina el rendimiento de estos modelos
en mayor medida que la elección del algoritmo. Se distinguen cuatro bloques:

\begin{itemize}
  \item \textbf{Desfases} ($\text{lag}_k$). El valor de la serie $k$ instantes
        antes. Su selección debe fundamentarse en la estructura de autocorrelación
        y, sobre todo, respetar la restricción de disponibilidad: para un
        horizonte de $h$ periodos únicamente son admisibles los desfases de orden
        $k \geq h$.

  \item \textbf{Estadísticos móviles}. Media y desviación típica sobre ventanas
        precedentes. Recogen el nivel y la volatilidad recientes. La ventana debe
        cerrarse antes del instante de predicción.

  \item \textbf{Codificación cíclica}. Las variables de calendario son
        circulares: la hora 23 es contigua a la hora 0. Su representación como
        enteros introduce una discontinuidad artificial que se evita mediante
        pares de funciones trigonométricas.

  \item \textbf{Variables de calendario}. Indicadores de festividad, día
        laborable y sus entornos. Su relevancia en demanda eléctrica es
        considerable, como se cuantifica en el capítulo de resultados.
\end{itemize}


% ╔═══════════════════════════════════════════════════════════════════════════╗
% ║  SUSTITUYE LA SECCIÓN 3.4                                                 ║
% ╚═══════════════════════════════════════════════════════════════════════════╝

\section{El mercado eléctrico español y la tarifa PVPC}

\subsection*{El mercado diario}

La electricidad se negocia en España a través del Operador del Mercado Ibérico de
Energía \parencite{omie}. El mercado diario casa oferta y demanda para cada una de
las veinticuatro horas del día siguiente mediante un mecanismo marginalista: todas
las ofertas aceptadas perciben el precio de la última unidad necesaria para cubrir
la demanda prevista.

De este mecanismo se deriva una consecuencia directa para el presente trabajo:
\textbf{el precio horario guarda una relación estrecha con la demanda prevista}.
Las horas de mayor consumo exigen recurrir a tecnologías de coste marginal
superior, lo que eleva el precio de casación. Predecir la demanda equivale, por
tanto, a anticipar de forma aproximada el perfil de precios.

\subsection*{La tarifa PVPC}

El Precio Voluntario para el Pequeño Consumidor es la tarifa regulada aplicable a
suministros de baja tensión con potencia contratada inferior a diez kilovatios.
Su característica distintiva es que el término de energía \textbf{varía cada hora}
en función del resultado del mercado mayorista.

Los precios correspondientes a cada día se publican la víspera, en torno a las
20:15 horas, una vez cerrada la casación del mercado diario. El consumidor acogido
a esta modalidad dispone así de información con la que planificar su consumo, si
bien limitada al día siguiente.

\subsection*{Implicaciones para el sistema propuesto}

De lo anterior se derivan las dos aportaciones que la aplicación desarrollada
puede ofrecer sobre la información ya disponible:

\begin{enumerate}
  \item \textbf{Anticipación.} La previsión de demanda está disponible en
        cualquier momento, mientras que el precio oficial solo se publica a las
        20:15 horas de la víspera.

  \item \textbf{Horizonte extendido.} El mercado publica únicamente el día
        siguiente; el sistema propuesto alcanza las cuarenta y ocho horas.
\end{enumerate}

Debe precisarse que la aplicación no predice el precio, sino la demanda, y emplea
esta última como indicador aproximado de aquel. La relación entre ambas magnitudes
es estrecha pero no determinista: la composición del parque generador disponible
—singularmente la aportación renovable— introduce una variabilidad que el modelo
no contempla. Esta limitación se recoge de forma explícita en la aplicación.


% ╔═══════════════════════════════════════════════════════════════════════════╗
% ║  SUSTITUYE LA SECCIÓN 3.5                                                 ║
% ╚═══════════════════════════════════════════════════════════════════════════╝

\section{Trabajos previos relacionados}

% ─────────────────────────────────────────────────────────────────────────────
% PENDIENTE: incorporar las referencias del director del trabajo,
% Gonzalo Surribas Sayago. Solicitar las citas completas y añadirlas
% a referencias.bib antes de la entrega.
% ─────────────────────────────────────────────────────────────────────────────

La revisión de \textcite{hong2016probabilistic} constituye la referencia de
síntesis más completa sobre predicción de carga eléctrica. Los autores clasifican
las aproximaciones existentes, discuten las métricas de evaluación habituales y
subrayan una carencia recurrente en la literatura: la ausencia de comparación
frente a líneas base ingenuas, lo que dificulta valorar la aportación real de las
propuestas. El presente trabajo atiende expresamente esa observación, evaluando
tres modelos ingenuos antes de cualquier formulación estadística.

En cuanto a la evidencia comparativa, las competiciones M constituyen el
experimento controlado de mayor alcance en el campo. La cuarta edición
\parencite{makridakis2020m4} estableció la superioridad de las combinaciones de
métodos y del aprendizaje transversal entre series. La quinta
\parencite{makridakis2022m5} confirmó el predominio de los métodos de potenciación
del gradiente sobre las arquitecturas de aprendizaje profundo, resultado
especialmente relevante para la orientación metodológica adoptada aquí.

Sobre los aspectos de descomposición estacional, el procedimiento STL
\parencite{cleveland1990stl} constituye la referencia clásica, y su extensión MSTL
\parencite{bandara2021mstl} permite el tratamiento de estacionalidades múltiples,
circunstancia que concurre en las series horarias de demanda.

Finalmente, en el terreno de la interpretabilidad, los valores SHAP
\parencite{lundberg2017shap} han desplazado a las medidas de importancia nativas
de los modelos de árboles, al proporcionar una descomposición aditiva de cada
predicción individual con fundamento en la teoría de juegos cooperativos.
